





\special{papersize=8.5in,11in}



\section{Artifact Appendix}

\subsection{Abstract}

 
This artifact contains necessary scripts and dependencies
to faithfully reproduce the crucial experiments presented in the paper. To successfully run the experiments, a host system with x86\_64 CPUs is required, along with at least one A100 or L40S NVIDIA GPU. We also provide a pre-built docker image to simplify the environment setup process. 


\subsection{Artifact check-list (meta-information)}


{\small
\begin{itemize}
  \item {\bf Program:} Efficiency benchmarking code for QServe; efficiency benchmarking code for baseline systems such as TensorRT-LLM.
  \item {\bf Compilation:} Completed in the docker.
  \item {\bf Transformations:} N/A. 
  \item {\bf Binary:} N/A.
  \item {\bf Model:} Llama-3-8B, Llama-2-7B, Mistral-7B, Llama-2-13B.
  \item {\bf Data set:} None.
  \item {\bf Run-time environment:} NVIDIA Container Toolkit (nvidiadocker).
  \item {\bf Hardware:} A host with x86\_64 CPUs and at least one NVIDIA A100 GPU (recommended) or L40S GPU.
  \item {\bf Run-time state:} N/A.
  \item {\bf Execution:} All benchmarks are executed on NVIDIA GPUs, while some data pre-processing code is executed on the host CPU.
  \item {\bf Metrics:} LLM generation throughput.
  \item {\bf Output:} Generation throughput (tokens/second).
  \item {\bf Experiments:} Inference speed measurement for QServe and baseline systems such as TensorRT-LLM.
  \item {\bf How much disk space required (approximately)?:} 512G.
  \item {\bf How much time is needed to prepare workflow (approximately)?:} Around 1 hour to pull docker images depending on the Internet connection and CPU performance.
  \item {\bf How much time is needed to complete experiments (approximately)?:} Around 1 hour to finish the efficiency benchmarks of QServe; and 2-4 GPU hours to finish the TensorRT-LLM benchmarks depending on the GPU performance and number of tasks to evaluate.
  \item {\bf Publicly available?:} Yes.
  \item {\bf Code licenses (if publicly available)?:} Apache License 2.0.
  \item {\bf Data licenses (if publicly available)?:} MIT.
  \item {\bf Workflow framework used?:} Docker.
  \item {\bf Archived (provide DOI)?:} \url{https://doi.org/10.5281/zenodo.14991385}
\end{itemize}

\subsection{Description}

\vspace{-3pt}
\subsubsection{How delivered}
\vspace{-3pt}

We will provide AE reviewers with a pre-built docker image containing \system, TensorRT-LLM and all necessary dependencies.

\vspace{-5pt}
\subsubsection{Hardware dependencies}
\vspace{-3pt}

A host machine with x86\_64 CPUs and at least one NVIDIA A100 GPU (recommended) or L40S GPU.

\vspace{-5pt}
\subsubsection{Software dependencies}
\vspace{-3pt}

A GPU-compatible Docker runtime environment is required.


\vspace{-5pt}
\subsection{Installation}

We recommend that users utilize our pre-built Docker images to set up the environment and run all experiments within the GPU-supported Docker container.



\begin{lstlisting}[style=mystyle, language=bash]
docker run --gpus all -it --workdir /root shang12138/qserve-mlsys25-ae
\end{lstlisting}

\vspace{-10pt}
\subsection{Experiment workflow}

The generation throughputs of QServe and baseline system (i.e., TensorRT-LLM) can be measured with the following commands.




\begin{lstlisting}[style=mystyle, language=bash]
# QServe benchmark
cd /root/QServe    
bash scripts/benchmark/benchmark_a100.sh 
# Run this command if you run on A100 GPU
# Results in ./parsed_results.csv

# TensorRT-LLM benchmark
cd /root/TensorRT-LLM   
bash launch-all.sh   
# Launch TensorRT-LLM evaluation
# Results in ./results.csv
\end{lstlisting}


\vspace{-10pt}
\subsection{Evaluation and expected result}
\vspace{-15pt}
\begin{table}[h]
\centering
\caption{Generation throughput of QServe and baseline (TensorRT-LLM). Unit: tokens/second.}
\begin{tabular}{cccc}
\toprule
Model       & TensorRT-LLM (W8A8KV8) & QServe     \\ \midrule
Llama-3-8B  & 2387.55  & 2980.69    \\
Llama-2-7B  & 2339.97  & 2860.01    \\
Mistral-7B  & 2427.64  & 3031.93    \\ \bottomrule
\end{tabular}
\label{appendix:expected_results}
\vspace{-5pt}
\end{table}

Table~\ref{appendix:expected_results} provides reference numbers for \system benchmarks. Please note that absolute throughput measurements may vary slightly, even on identical GPU platforms, due to differences in machine conditions. However, the relative acceleration ratios should remain consistent.



\subsection{Experiment customization}

The users are encouraged to carry out experiments with different models and batch sizes by modifying the benchmarking scripts.
Accuracy evaluation is omitted to simplify the environment setup. The accuracy results can be reproduced with open-source library \href{https://github.com/mit-han-lab/deepcompressor/tree/dev/v0.1.0}{\texttt{deepcompressor}}.


\subsection{Methodology}

Submission, reviewing and badging methodology:
\vspace{-10pt}
{\small
\begin{itemize}
  \item \url{http://cTuning.org/ae/submission-20190109.html}
  \vspace{-5pt}
  \item \url{http://cTuning.org/ae/reviewing-20190109.html}
  \vspace{-5pt}
  \item \url{https://www.acm.org/publications/policies/artifact-review-badging}
\end{itemize}
}


